\documentclass{article}
\usepackage{../math_template}

\author{Jonathan Kasongo}
\title{JEE Mains 2025 Shift 2 --- Q6}

\begin{document}
\maketitle
\textit{An educated guess is not a guess.}
\begin{problem}{JEE Mains 2025 Shift 2 --- Q6}
If
$$
\int e^x \left( \frac{x\arcsin x}{\sqrt{1-x^2}} +
                \frac{\arcsin x}{(1-x^2)^{\frac32}} +
                \frac{x}{1-x^2}
\right) \mathrm{d}x = g(x) + C
$$
where $C$ is the constant of integration, then $g(1/2)$ is?
\end{problem}

\begin{solution}{Solution}
It is natural to guess that the anti-derivative might take the form
$g(x) = e^x q(x) \arcsin x$. We have
$$
g'(x) = e^x q(x) \arcsin x + e^x q'(x)\arcsin x +
e^x q(x)\frac{1}{\sqrt{1-x^2}}.
$$
We want to try and match that last term to the last term in the integrand
so we try the function
$$
q(x) = \frac{x}{\sqrt{1-x^2}}
$$
which has derivative
$$
q'(x) = \frac{1}{\sqrt{1-x^2}} + \frac{x^2}{(1-x^2)^{\frac32}}.
$$
So we can re-write $g'(x)$ now as
\[
\begin{aligned}
g'(x) &= e^x \left( \frac{x \arcsin x}{\sqrt{1-x^2}} +
                    \frac{\arcsin x}{\sqrt{1-x^2}}   +
                    \frac{x^2 \arcsin x}{(1-x^2)^\frac32} +
                    \frac{x}{1-x^2}
                    \right)\\
&= e^x \left( \frac{x \arcsin x}{\sqrt{1-x^2}} +
              \frac{(1-x^2)\arcsin (x) + x^2\arcsin x}{(1-x^2)^{\frac32}} +
              \frac{x}{1-x^2}
\right) \\
&= e^x \left( \frac{x\arcsin x}{\sqrt{1-x^2}} +
              \frac{\arcsin x}{(1-x^2)^{\frac32}} +
              \frac{x}{1-x^2}
\right)\\
\end{aligned}
\]
which is the given integrand. So we know that our choice of $g(x)$ does
in fact satisfy the given equation. Finally we can just compute
$$
g\left(\frac12\right) = \frac{\pi}{6}\sqrt{\frac{e}{3}}. \ \blacksquare
$$

\textit{Comments: Interestingly enough ChatGPT with thinking mode,
struggled significantly to solve this one even after I gave hints. The
real trap is that after you compute $g'(x)$ it doesn't look like it
matches the integrand, but if you combine like terms correctly, you'll
see that it actually does match the integrand.}
\end{solution}

\end{document}
